% GENERATED FILE. Edit canonical lessons, then run npm run generate:books.
% canonical-source-hash: fnv1a64:6cd4fe14bd21dd5f
% canonical-lessons: ML-C42-big, ML-C42-small, ML-C42-good, ML-C42-new, ML-C42-old, ML-C42-adjective-system

\chapter{Describing Things}
\label{ch:ml-describing-things}
\hlchaptermodality{42}

\begin{chapteropening}
\textbf{By the end of this chapter:} \emph{I can point at something in Malayalam and ask who or where it is — and I can say what the six words have in common.}

Say the near word, the far word and the question word in a row, and name the part that changed.
\end{chapteropening}

\section[valiya]{\ml{വലിയ} (valiya) — big}

\label{lesson:ML-C42-big}



\begin{quote}
You can already name things. This is how you POINT at them.
\end{quote}

\subsection*{You'll want to know: \ml{വലിയ}}
\textbf{\ml{വലിയ}} — \emph{valiya} — big.

Put it in front of a word you already know and you have said something new. That is all an adjective does in this language.

The last lesson of this chapter gives the one rule that governs every adjective you will ever meet here.

\subsection*{Your turn}
Say these aloud:

\begin{itemize}
\raggedright
  \item \textquotedblleft{}\ml{വലിയ}\textquotedblright{} three times, pointing at something different each time
  \item it once with a word you already know after it
\end{itemize}

\begin{tcolorbox}[breakable,colback=teal!4,colframe=teal!35!black,title={Before you move on}]
What is \emph{valiya}? (\textbf{big}.) Where does it go — before the word it describes, or after? (\textbf{Before.})
\end{tcolorbox}
\section[ceṟiya]{\ml{ചെറിയ} (ceṟiya) — small}

\label{lesson:ML-C42-small}



\begin{quote}
You can already name things. This is how you POINT at them.
\end{quote}

\subsection*{You'll want to know: \ml{ചെറിയ}}
\textbf{\ml{ചെറിയ}} — \emph{ceṟiya} — small.

Put it in front of a word you already know and you have said something new. That is all an adjective does in this language.

You met \textbf{\ml{വലിയ}} \emph{valiya} a moment ago. Put them side by side and nothing about either one changes.

The last lesson of this chapter gives the one rule that governs every adjective you will ever meet here.

\subsection*{Your turn}
Say these aloud:

\begin{itemize}
\raggedright
  \item \textquotedblleft{}\ml{ചെറിയ}\textquotedblright{} three times, pointing at something different each time
  \item it once with a word you already know after it
\end{itemize}

\begin{tcolorbox}[breakable,colback=teal!4,colframe=teal!35!black,title={Before you move on}]
What is \emph{ceṟiya}? (\textbf{small}.) Where does it go — before the word it describes, or after? (\textbf{Before.})
\end{tcolorbox}
\section[nalla]{\ml{നല്ല} (nalla) — good}

\label{lesson:ML-C42-good}



\begin{quote}
You can already name things. This is how you POINT at them.
\end{quote}

\subsection*{You'll want to know: \ml{നല്ല}}
\textbf{\ml{നല്ല}} — \emph{nalla} — good.

Put it in front of a word you already know and you have said something new. That is all an adjective does in this language.

You met \textbf{\ml{ചെറിയ}} \emph{ceṟiya} a moment ago. Put them side by side and nothing about either one changes.

The last lesson of this chapter gives the one rule that governs every adjective you will ever meet here.

\subsection*{Your turn}
Say these aloud:

\begin{itemize}
\raggedright
  \item \textquotedblleft{}\ml{നല്ല}\textquotedblright{} three times, pointing at something different each time
  \item it once with a word you already know after it
\end{itemize}

\begin{tcolorbox}[breakable,colback=teal!4,colframe=teal!35!black,title={Before you move on}]
What is \emph{nalla}? (\textbf{good}.) Where does it go — before the word it describes, or after? (\textbf{Before.})
\end{tcolorbox}
\section[putiya]{\ml{പുതിയ} (putiya) — new}

\label{lesson:ML-C42-new}



\begin{quote}
You can already name things. This is how you POINT at them.
\end{quote}

\subsection*{You'll want to know: \ml{പുതിയ}}
\textbf{\ml{പുതിയ}} — \emph{putiya} — new.

Put it in front of a word you already know and you have said something new. That is all an adjective does in this language.

You met \textbf{\ml{നല്ല}} \emph{nalla} a moment ago. Put them side by side and nothing about either one changes.

The last lesson of this chapter gives the one rule that governs every adjective you will ever meet here.

\subsection*{Your turn}
Say these aloud:

\begin{itemize}
\raggedright
  \item \textquotedblleft{}\ml{പുതിയ}\textquotedblright{} three times, pointing at something different each time
  \item it once with a word you already know after it
\end{itemize}

\begin{tcolorbox}[breakable,colback=teal!4,colframe=teal!35!black,title={Before you move on}]
What is \emph{putiya}? (\textbf{new}.) Where does it go — before the word it describes, or after? (\textbf{Before.})
\end{tcolorbox}
\section[paḻaya]{\ml{പഴയ} (paḻaya) — old — of a thing, not of a person}

\label{lesson:ML-C42-old}



\begin{quote}
You can already name things. This is how you POINT at them.
\end{quote}

\subsection*{You'll want to know: \ml{പഴയ}}
\textbf{\ml{പഴയ}} — \emph{paḻaya} — old — of a thing, not of a person.

Put it in front of a word you already know and you have said something new. That is all an adjective does in this language.

You met \textbf{\ml{പുതിയ}} \emph{putiya} a moment ago. Put them side by side and nothing about either one changes.

The last lesson of this chapter gives the one rule that governs every adjective you will ever meet here.

\subsection*{Your turn}
Say these aloud:

\begin{itemize}
\raggedright
  \item \textquotedblleft{}\ml{പഴയ}\textquotedblright{} three times, pointing at something different each time
  \item it once with a word you already know after it
\end{itemize}

\begin{tcolorbox}[breakable,colback=teal!4,colframe=teal!35!black,title={Before you move on}]
What is \emph{paḻaya}? (\textbf{old — of a thing, not of a person}.) Where does it go — before the word it describes, or after? (\textbf{Before.})
\end{tcolorbox}
\section[before the noun, unchanged]{before the noun, unchanged — five words, one rule}

\label{lesson:ML-C42-adjective-system}



\begin{quote}
Say the five words from this chapter out loud, then say each one in front of a noun you already know.
\end{quote}

\begin{grammarlens}[title={Grammar lens: where an adjective goes}]
You did not learn five separate words. You learned \textbf{where an adjective goes}, and five words that go there.

\noindent
\begin{tabularx}{\linewidth}{@{}>{\raggedright\arraybackslash}X>{\raggedright\arraybackslash}X@{}}
\toprule
\textbf{meaning} & \textbf{word} \\
\midrule
big & \textbf{\ml{വലിയ}} \emph{valiya} \\
small & \textbf{\ml{ചെറിയ}} \emph{ceṟiya} \\
good & \textbf{\ml{നല്ല}} \emph{nalla} \\
\bottomrule
\end{tabularx}

Every one ends in -a, sits in front of the noun, and never changes.

This is worth more than the five words are. Every noun you already know — your name, the house, the water, the day — can now take any of them in front, and every adjective you meet from here goes in the same place.
\end{grammarlens}

\subsection*{Your turn}
\begin{itemize}
\raggedright
  \item \emph{Say it:} each of the five in front of a noun you already know
  \item \emph{Say it:} two of them with the same noun, one after the other
  \item \emph{Notice:} whether the adjective changed shape when the noun did
\end{itemize}

\begin{tcolorbox}[breakable,colback=teal!4,colframe=teal!35!black,title={Before you move on}]
Where does an adjective go in this language? (\textbf{In front of the noun.}) Does it change shape to match the noun? (unchanged.)
\end{tcolorbox}
